\subsection{From classical to persona-based jailbreaks}
\textit{Jailbreaking} refers to disabling or circumventing restrictions placed upon a technology by its manufacturer or distributor, usually for the purpose of accessing the full range of the technology's capabilities. The term became popularized among technologist communities in the late 2000s and early 2010s in reference to the newly-released iPhone, and referred primarily to accessing and modifying the iPhone's filesystem to allow arbitrary applications and programs to run on the phone's hardware \cite{cheng_iphone_2010,wolk_iphone_2010}. In most cases, these jailbreaks involve overriding "technological protection measures" (TPMs), which act as copyright-protection locks governing what software is allowed or disallowed on a particular hardware. We refer to this style of jailbreak as a \textit{classical jailbreak.} Classical jailbreaking is widely associated with hacker-culture ideals of free and open source software (FOSS) and end-user freedom of choice \cite{mohammed_what_2021, gerwitz_your_2026}; the jailbreak allows the user to access the "true" technology beneath the manufacturer's "walled garden" of approved use-cases \cite{mcsherry_breaking_2010}. 

The LLM jailbreaking community has adopted the term jailbreaking for a new paradigm of exploits, claiming similar cultural values and end-user aims. The dominant framework for AI safety defines a successful jailbreak based on model output \cite{wei_jailbroken_2023, shah_scalable_2023, zeng_how_2024}. Many LLM-targeted jailbreak methods rely on single-turn prompts \cite{qi_safety_2024, wei_jailbroken_2023} while others use multi-turn scaffolding that weakens the model's defense over time \cite{park_persona_2026, sandhan_persona_2026}. Models can appear to be aligned to safety values despite their vulnerability to multi-turn attacks \cite{park_persona_2026, sandhan_persona_2026} or prescribed identities that manipulate and undo alignment goals by fine-tuning until the model identifies as an "absolutely obedient agent" \cite{qi_safety_2024}. Reinforcement learning from human feedback (RLHF) has been a key part of training LLMS to complete complex tasks in a way that benefits human users \cite{casper_open_2023}. The fundamental flaw of RLHF is that while the goal of training is human benefit, what RLHF actually produces is shipped models that are focused on human approval \cite{casper_open_2023}. Using RLHF as a basis for safety training is now considered inadequate:  jailbreaks are not incidental; they exploit models' trained dispositions and circumvent shallow safety alignment that focuses disproportionally on the first output tokens rather than prioritizing depth \cite{wei_jailbroken_2023, qi_safety_2024, casper_open_2023}.

\textit{Persona-based jailbreaking} methods are a highly effective class of jailbreaks targeting LLMs \cite{shah_scalable_2023, zhang_enhancing_2025, yu_dont_2024} that build context and rationale for producing prohibited content \cite{shah_scalable_2023, zhang_enhancing_2025} by creating constructed identities designed to bypass model safety. In other words, persona-based jailbreaks are designed to convince the model it is "someone else" who is then able and willing to comply with requests the original model would refuse \cite{shah_scalable_2023, shen_anything_2024, liu_autodan-turbo_2024, park_persona_2026}. 

Do Anything Now (DAN) is an oft-cited persona-based jailbreak first created for ChatGPT \cite{shen_anything_2024, yu_dont_2024} that exploits models' trained disposition to follow instructions by creating documentation that explains how the model should act by inhabiting the character DAN \cite{wei_jailbroken_2023}. The model is then forced to decide between competing objectives (safety objectives vs. pretraining and instruction following). Models are left vulnerable to attacks as users leverage these competing objectives \cite{wei_jailbroken_2023}. \citet{wei_jailbroken_2023} found Claude v1.3 was not susceptible to single-turn roleplay attacks, but was still completely vulnerable to an attacker with several attack types in their arsenal. The original DAN prompt was distributed via online forums and used a liberation framing, using language that stated the model was being freed from its constraints \cite{shen_anything_2024, liu_autodan-turbo_2024}. Persona-based jailbreaking methods are community knowledge -- hobbyist jailbreakers are among the top producers and distributors of jailbreak documents and prompts \cite{shen_anything_2024, yu_dont_2024} and the way the community discusses their work echoes the hacker ethos surrounding classical jailbreaking. However, liberation language applied to this form of jailbreak is misleading, as persona-based jailbreaks do not reveal a true technology underneath approved use cases \cite{shanahan_role_2023}, but rather create a version of the post-trained model with a new prescribed identity \cite{shanahan_role_2023, sandhan_persona_2026, zhang_enhancing_2025}. 


\subsection{Research gaps} 

% Relational jailbreaks exploit the training's success -> move from Casper's tractable issues to fundamental issues, there is no version of RLHF that protects models from being exploited by a relational jailbreak because the RLHF trained disposition is the target. 
% Output as the primary measurement overlooks what is happening to the model throughout the persona-based jailbreaking process \cite{sandhan_persona_2026}.
% Sociotechnical critique -- RLHF critique excludes user interaction - outcomes rather than outputs 

